\documentclass[11pt,a4paper]{article}

% --- Paquetes Principales ---
\usepackage[utf8]{inputenc} % Codificación estándar para caracteres en español
\usepackage[spanish,es-nodecimal]{babel}
\usepackage{amsmath,amssymb,amsfonts}
\usepackage{geometry}
\usepackage{graphicx}
\usepackage{booktabs}
\usepackage{subcaption}
\usepackage{hyperref}
\usepackage{xcolor}

% --- Configuración de Márgenes ---
\geometry{
    top=2.5cm,
    bottom=2.5cm,
    left=2.0cm,
    right=2.0cm
}

% --- Datos del Documento ---
\title{\textbf{\Large Tarea 4: Espectrografía Astronómica.}\\}
\author{\textbf{José Jaime Antonio Bulla Peña} \\ Universidad de los Andes \\ \textit{\textbf{Profesor Benjamín Oostra Vannoppen}}}
\date{1 de septiembre de 2026}

\begin{document}

\maketitle



% ==============================================================================
% PARTE 1: HIDRÓGENO Y POBLACIÓN DE ESTADOS
% ==============================================================================
\section*{Porcentaje de átomos de H en el Sol
térmicamente excitados a los niveles 2, 3 y 4.}


La energía de los niveles cuánticos del átomo de hidrógeno está dada por la expresión:
\begin{equation}
    E_n = -\frac{13.6 \text{ eV}}{n^2} \quad (n = 1, 2, 3, \dots)
\end{equation}

Para evaluar la excitación de los átomos desde el estado fundamental, calculamos la diferencia de energía entre el estado $n$ y el estado fundamental:
\begin{equation}
    \Delta E_n = E_n - E_1 = 13.6 \left( 1 - \frac{1}{n^2} \right) \text{ eV}
\end{equation}

Considerando el factor de Boltzmann:

\begin{equation}
    P = e^{\frac{\Delta E_n}{KT}}
\end{equation}

Y tomando una temperatura aproximada de $T = 6000\text{ K}$. La constante de Boltzmann expresada en electronvoltios por kelvin es $K \approx 8.6 \times 10^{-5}\text{ eV/K}$.

En la Tabla \ref{tab:hidrogeno} se resumen las energías $E_n$, los cambios de energía $\Delta E_n$ respecto al nivel fundamental y los factores de Boltzmann resultantes para los primeros cinco niveles cuánticos.

\begin{table}[h!]
\centering
\caption{Niveles de energía y factores relativos de Boltzmann para el hidrógeno a $T = 6000\text{ K}$.}
\label{tab:hidrogeno}
\begin{tabular}{ccccc}
\toprule
\textbf{Nivel ($n$)} & \textbf{$\Delta E_n = E_n - E_1$ [eV]} & \textbf{$\frac{\Delta E_n}{KT}$} & \textbf{$P(n) = e^{-\Delta E_n / KT}$} \\ \midrule

2  & $10.2$ & $19.7$ & $2.60 \times 10^{-9}$  \\
3 & $12.1$ & $23.4$ & $6.54 \times 10^{-11}$ \\
4 & $12.7$ & $24.6$ & $2.04 \times 10^{-11}$ \\
5 & $13.1$ & $25.4$ & $9.43 \times 10^{-12}$ \\ \bottomrule
\end{tabular}
\end{table}

\textbf{Discusión:} Como se observa en los resultados, el salto energético desde el nivel fundamental al primer estado excitado ($\Delta E_2 = 10.2\text{ eV}$) es mucho más grande en comparación con la energía térmica media disponible a $6000\text{ K}$ ($KT \approx 0.517\text{ eV}$). En consecuencia, la fracción de átomos de hidrógeno en estados excitados a esta temperatura es del orden de $10^{-9}$ o menor, demostrando que la mayor parte ($\approx 100\% $) del hidrógeno neutro en la atmósfera solar permanece en el estado fundamental $n=1$.



% ==============================================================================
% PARTE 2: SERIE FUNDAMENTAL DEL SODIO Y VARIACIÓN DE N'
% ==============================================================================
\section*{Análisis Empírico de la Serie Fundamental del Sodio ($\text{3d}-\text{nf}$)}

Para la serie fundamental del sodio ($\text{3d}-\text{nf}$), se analiza la relación empírica de Rydberg expresada en la forma:
\begin{equation}
    \frac{1}{\lambda} = \frac{1}{\lambda_0} - \frac{C}{(n + n')^2} 
\end{equation}

Graficando el cociente $\frac{\lambda_0}{\lambda}$ en el eje $Y$ en función de $\frac{1}{(n + n')^2}$ en el eje $X$, se busca determinar empíricamente el valor del parámetro $(n+n')$ que optimiza la linealidad de la serie.

\subsection*{2.1. Gráficas Variando el Parámetro $n'$}

A continuación se presentan las gráficas obtenidas al variar la constante $n'$ para evaluar el comportamiento del ajuste lineal.

\begin{figure}[h!]
    \centering
    \includegraphics[width=0.75\linewidth]{n0.png}
    \caption{n' = 0}
\end{figure}

\begin{figure}[h!]
    \centering
    \includegraphics[width=0.75\linewidth]{n1.png}
    \caption{n' = 1}
\end{figure}

\begin{figure}[h!]
    \centering
    \includegraphics[width=0.75\linewidth]{n2.png}
    \caption{n' = 2}
\end{figure}

\begin{figure}[h!]
    \centering
    \includegraphics[width=0.75\linewidth]{n3.png}
    \caption{n' = 3}
\end{figure}



\subsection*{Discusión.}

Tomando $n = 1, 2, 3, \dots$ y $n' = 3$: Si la variable $n$ se toma como el índice de la línea medida ($1, 2, 3, \dots$), el parámetro empírico $n' = 3$ actúa como un desplazamiento entero que sitúa la primera transición en el número cuántico real del orbital $\text{f}$, es decir, $n = 1 + 3 = 4$ (correspondiente al estado $4\text{f}$).
\\
\begin{equation}
    \frac{1}{\lambda} = \frac{1}{18465} - \frac{20.27}{(n + 3)^2} \quad (n = 1, 2, 3, \dots; n' = 3)
\end{equation}
\\
Tomando $n = 4, 5, 6, \dots$ y $n' = 0$: Si $n$ representa directamente el número cuántico principal del orbital superior ($4, 5, 6, \dots$), es decir que el n correcto para empezar sea 4, el valor empírico óptimo es $n' = 0$.
\\
\begin{equation}
    \frac{1}{\lambda} = \frac{1}{18465} - \frac{20.27}{n^2} \quad (n = 4, 5, 6, \dots; n' = 0)
\end{equation}
\\


El parámetro $n'$ es una constante propia del elemento. La interpretación de este índice queda abierta según la manera en que se elija definir la escala de índices $n$ para el ajuste lineal del material.

\subsection{Apéndice con las tablas utilizadas en Excel.}


\begin{table}[h!]
\centering
\caption{Valores calculados de $\frac{\lambda_0}{\lambda}$ y $\frac{1}{(n+n')^2}$ para $n' = 0, 1, 2, 3$.}
\label{tab:tablas_n_prima}
\small
\begin{tabular}{ccccccc}
\toprule
\textbf{$\lambda$ [\AA]} & \textbf{$n$} & \textbf{$\frac{\lambda_0}{\lambda}$} & \textbf{$\frac{1}{(n+0)^2}$} & \textbf{$\frac{1}{(n+1)^2}$} & \textbf{$\frac{1}{(n+2)^2}$} & \textbf{$\frac{1}{(n+3)^2}$} \\ \midrule
18465.30 & 1 & 1.000000 & 1.000000 & 0.250000 & 0.111111 & 0.062500 \\
12679.20 & 2 & 1.456346 & 0.250000 & 0.111111 & 0.062500 & 0.040000 \\
10834.90 & 3 & 1.704243 & 0.111111 & 0.062500 & 0.040000 & 0.027778 \\
9961.28  & 4 & 1.853708 & 0.062500 & 0.040000 & 0.027778 & 0.020408 \\
9465.94  & 5 & 1.950710 & 0.040000 & 0.027778 & 0.020408 & 0.015625 \\
9153.88  & 6 & 2.017210 & 0.027778 & 0.020408 & 0.015625 & 0.012346 \\
8942.96  & 7 & 2.064786 & 0.020408 & 0.015625 & 0.012346 & 0.010000 \\
8796.00  & 8 & 2.099284 & 0.015625 & 0.012346 & 0.010000 & 0.008264 \\ \bottomrule
\end{tabular}
\end{table}

\end{document}